\input{ieee-preamble}
\hypersetup{
  pdftitle={Underactuated Humanoid Locomotion},
  pdfauthor={Santiago Restrepo},
  pdfsubject={Deep reinforcement learning method for humanoid bipedal walking},
  pdfkeywords={humanoid, bipedal locomotion, underactuation, deep learning, PPO, sim-to-real}
}

\title{Underactuated Humanoid Locomotion}
\author{\IEEEauthorblockN{Santiago Restrepo}
\IEEEauthorblockA{\href{mailto:contact@waajacu.com}{contact@waajacu.com}\quad
\url{https://waajacu.com/}\\[3pt]
\textit{Learning-based control method for bipedal walking}}}
\date{}

\begin{document}
\maketitle

\begin{abstract}
This brief proposes deep reinforcement learning for underactuated humanoid
walking. A neural policy learns foot placement, load transfer, and whole-body
coordination through interaction with simulated contact dynamics. Joint
feedback executes the learned commands. The workflow combines physical
modeling, policy training, and validation before hardware deployment; no
robot-specific performance is assumed.
\end{abstract}

\begin{IEEEkeywords}
Humanoid, bipedal locomotion, underactuation, deep reinforcement learning.
\end{IEEEkeywords}

\balance
\section{Model the Floating Base and Contact}
Assume a free-standing humanoid with $n$ actuated joints. Let $q$ denote
configuration and $v$ generalized velocity. Between impacts, the
floating-base dynamics are~\cite{tedrakeContact}
\begin{equation}
\begin{aligned}
M(q)\dot v+h(q,v)&=B\tau+J_c^{\mathsf T}\lambda,\\
\dot q&=N(q)v.
\end{aligned}
\end{equation}
$M$ is inertia, $h$ includes gravity and velocity terms, $B$ maps joint
torques $\tau$, and $J_c^{\mathsf T}$ maps contact forces $\lambda$; $N$
handles orientation kinematics. The base has six unactuated directions.
Foot contact changes the available control authority: fixed flat-foot
support can be fully actuated in reduced coordinates.

Model unilateral, nonpenetrating contact, friction limits, and impact
velocity resets. Foot forces come from the coupled contact dynamics. The
policy influences base momentum through joint motion and ground reaction
forces; it must learn when to shift weight and place the next foot.

\section{Connect a Neural Policy to the Motors}
Use command $c$ = (forward speed, lateral speed, yaw rate)
and history $H$ of IMU angular velocity, projected gravity, joint
positions/velocities, and previous actions. A causal Transformer encodes
this history to address partial observability~\cite{radosavovic2024}.
\begin{equation}
\begin{aligned}
a_t&\sim\pi_\theta(\,\cdot\mid H_t,c_t),\\
q_{\mathrm{des}}&=q_{\mathrm{nom}}+D\tanh(a_t),\\
\tau&=\operatorname{sat}[K_p(q_{\mathrm{des}}-q_a)-K_dv_a].
\end{aligned}
\end{equation}
$\theta$ denotes policy parameters; $D$ scales offsets about
$q_{\mathrm{nom}}$, and $\operatorname{sat}$ applies motor torque limits.
Joint states $q_a$, $v_a$ drive faster proportional-derivative (PD) loops.
Match gains, action rates, and delays to the hardware.

An asymmetric training critic can use simulator ground truth. The actor
must use only deployment observations; privileged terrain or velocity
signals must not leak into its inputs.

\section{Train for Commanded Locomotion}
Train an actor and value critic with proximal policy optimization
(PPO)~\cite{schulman2017} in parallel physics simulations. Collect rollouts,
estimate advantages, update the actor with PPO's clipped objective, and
fit the critic to return targets. Optimize the expected discounted return:
\begin{equation}
\max_\theta J(\theta)=\mathbb E_{\pi_\theta}
\left[\sum_{t=0}^{T-1}\gamma^t r_t\right].
\end{equation}
$T$ is the task horizon and $\gamma$ the discount factor; bootstrap values
at nonterminal rollout truncations. Reward velocity/yaw tracking and
upright locomotion; penalize falls, slip, collisions, effort, and abrupt
action changes. Normalize terms before tuning weights. Test for shortcuts
such as shuffling or hopping when the task calls for walking.

\section{Bridge Simulation and Hardware}
Identify link inertia, joint friction, motor response, and latency.
Randomize mass, center of mass, ground friction, actuator strength, sensor
noise, and delays over plausible ranges. Add pushes and a curriculum from
standing and slow steps to turns and uneven terrain. Such
simulation-trained policies have transferred to real
humanoids~\cite{radosavovic2024}; transfer must be verified for the target
robot.

Freeze network weights and observation normalization for deployment. Use
the policy mean with the same action transformation and filtering.
Enforce actuator and command limits, monitor joint state and temperature,
and provide an independent timeout/fall supervisor.

\section{Measure What Generalizes}
Evaluate multiple training seeds on held-out terrains, payloads,
disturbances, and latency. Compare with a model-based walking baseline
using fall rate, command-tracking error, slip, energy per distance, torque
saturation, and inference time. Progress from simulation to restrained
hardware trials and supervised walking. Report tested conditions and
failures; a successful rollout is not a stability proof.

\begingroup
\raggedright
\bibliographystyle{IEEEtran}
\bibliography{references}
\endgroup
\end{document}
